% ============================================================================
%  ESTADÍSTICA INFERENCIAL — Sesión 13
%  Semana 13: 2 al 7 de noviembre de 2026   (lunes 2: festivo)
%  3.1 Concepto de pruebas de hipótesis · 3.2 Tipos de pruebas de hipótesis
% ============================================================================
\documentclass[aspectratio=169,10pt]{beamer}

\usepackage[utf8]{inputenc}
\IfFileExists{spanish.ldf}{\usepackage[spanish,es-noshorthands]{babel}}{}
\usepackage{amsmath,amssymb}
\usepackage{booktabs}
\usepackage{colortbl}
\usepackage{tikz}
\usepackage{listings}

\usetheme{Uniminuto}

\uniasignatura{Estadística inferencial}
\uniprofesor{Andrés Felipe Puerta Velez}
\unicorreo{andres.puerta-v@uniminuto.edu.co}
\uninrc{CBASBL061}
\unisemana{Sesión 13 — Semana 13 · 2 al 7 de noviembre de 2026}

\newcommand{\listatiny}{\setbeamerfont{itemize/enumerate body}{size=\tiny}\tiny}

\newcommand{\fx}[1]{{\ttfamily\color{uniAzul}#1}}

% ------------------------------------------------------------- CÓDIGO R ----
\lstdefinestyle{Rcode}{
  language=R,
  basicstyle=\ttfamily\scriptsize,
  keywordstyle=\color{uniAzul}\bfseries,
  commentstyle=\color{uniGris}\itshape,
  stringstyle=\color{uniTitulo},
  showstringspaces=false, numbers=none, breaklines=true,
  columns=fullflexible,
  backgroundcolor=\color{uniCremaClaro},
  frame=leftline, framerule=1.2pt, rulecolor=\color{uniNaranja},
  xleftmargin=0.35cm, framexleftmargin=0.35cm,
  framexrightmargin=0.15cm, framextopmargin=2pt, framexbottommargin=2pt,
  aboveskip=0.15cm, belowskip=0.15cm
}
\lstdefinestyle{Rmini}{
  language=R,
  basicstyle=\ttfamily\tiny,
  keywordstyle=\color{uniAzul}\bfseries,
  commentstyle=\color{uniGris}\itshape,
  stringstyle=\color{uniTitulo},
  showstringspaces=false, numbers=none, breaklines=true,
  columns=fullflexible,
  backgroundcolor=\color{uniCremaClaro},
  frame=leftline, framerule=1.2pt, rulecolor=\color{uniNaranja},
  xleftmargin=0.35cm, framexleftmargin=0.35cm,
  framexrightmargin=0.15cm, framextopmargin=2pt, framexbottommargin=2pt,
  aboveskip=0.12cm, belowskip=0.12cm
}
\lstdefinestyle{Rout}{
  basicstyle=\ttfamily\tiny\color{black},
  showstringspaces=false, numbers=none, breaklines=true,
  columns=fullflexible,
  backgroundcolor=\color{uniGris!10},
  frame=leftline, framerule=1.2pt, rulecolor=\color{uniGris},
  xleftmargin=0.35cm, framexleftmargin=0.35cm,
  framexrightmargin=0.15cm, framextopmargin=2pt, framexbottommargin=2pt,
  aboveskip=0.12cm, belowskip=0.12cm
}

\begin{document}

\portada

\begin{frame}{Orden del día}
  \begin{enumerate}
    \item Llamado a lista y retroalimentación del Examen 2
    \item Presentación de la Unidad 3
    \item \azul{Tema 3.1: concepto de prueba de hipótesis}
    \item \azul{Tema 3.2: tipos de pruebas de hipótesis}
    \item Errores tipo I y II, nivel de significancia y valor $p$
    \item Ejercicios en clase y cierre
  \end{enumerate}
\end{frame}

\begin{frame}{Dónde estamos: empieza la Unidad 3}
  En la unidad 2 terminamos comparando un intervalo con un valor de
  referencia: ``el intervalo $(0{,}5724;\ 0{,}6676)$ está por encima de
  $0{,}60$, luego se cumple la meta''.

  \vspace{0.2cm}
  \begin{alertblock}{La idea de la Unidad 3}
    \footnotesize
    Ese razonamiento informal se vuelve un \destacar{procedimiento formal de
    decisión}: la prueba de hipótesis. Permite responder con un criterio claro
    preguntas del tipo ``¿el nuevo proceso sí mejoró?'' o ``¿la máquina está
    descalibrada?''.
  \end{alertblock}

  \vspace{0.15cm}
  {\footnotesize\color{uniGris} Unidad 3: concepto y tipos de pruebas,
  planteamiento y pasos, software estadístico, correlación y regresión lineal.
  \destacar{Examen 3: semana 16}.}
\end{frame}

\seccion{3.1 Concepto de prueba de hipótesis}

\begin{frame}{Las dos hipótesis}
  \begin{columns}[T,onlytextwidth]
    \begin{column}{0.48\textwidth}
      \begin{block}{Hipótesis nula $H_0$}
        \footnotesize
        La afirmación que se somete a prueba. Representa el
        \azul{statu quo}: ``no hay cambio'', ``no hay diferencia'', ``se cumple
        el valor de referencia''.

        \vspace{0.1cm}
        Siempre lleva el signo de igualdad: $=$, $\le$ o $\ge$.
      \end{block}
    \end{column}
    \begin{column}{0.48\textwidth}
      \begin{exampleblock}{Hipótesis alterna $H_1$}
        \footnotesize
        Lo que el investigador quiere demostrar. Es la afirmación
        \azul{contraria} y nunca lleva igualdad: $\neq$, $>$ o $<$.
      \end{exampleblock}
    \end{column}
  \end{columns}

  \vspace{0.25cm}
  \begin{alertblock}{La lógica es la del juicio penal}
    \footnotesize
    $H_0$ es ``inocente'' y se mantiene mientras no haya evidencia suficiente.
    La muestra es la prueba. Si la evidencia es fuerte, se \destacar{rechaza}
    $H_0$; si no, se dice que \destacar{no hay evidencia suficiente para
    rechazarla}, que no es lo mismo que declararla verdadera.
  \end{alertblock}
\end{frame}

\begin{frame}{Cómo se plantean}
  \footnotesize
  \begin{center}
  \begin{tabular}{@{}p{5.4cm}p{2.6cm}p{2.6cm}@{}}
    \toprule
    \textbf{Afirmación del problema} & \textbf{$H_0$} & \textbf{$H_1$} \\
    \midrule
    ``El peso promedio es 2{,}5 kg'' & $\mu=2{,}5$ & $\mu\neq 2{,}5$ \\
    ``El nuevo proceso reduce el tiempo (antes 15 min)'' &
      $\mu\ge 15$ & $\mu<15$ \\
    ``Más del 60\% está satisfecho'' & $p\le 0{,}60$ & $p>0{,}60$ \\
    ``La máquina no está descalibrada'' & $\mu=\mu_0$ & $\mu\neq\mu_0$ \\
    \bottomrule
  \end{tabular}
  \end{center}

  \vspace{0.25cm}
  \begin{alertblock}{Regla para no equivocarse}
    \footnotesize
    Lo que se quiere \destacar{probar} va en $H_1$. Si el enunciado dice
    ``demuestre que'', ``verifique si mejora'', ``hay evidencia de que'', eso es
    $H_1$. La igualdad siempre queda en $H_0$.
  \end{alertblock}
\end{frame}

\seccion{3.2 Tipos de pruebas de hipótesis}

\begin{frame}{Unilateral o bilateral}
  \footnotesize
  \begin{center}
  \begin{tabular}{@{}llp{5.6cm}@{}}
    \toprule
    \textbf{Tipo} & \textbf{$H_1$} & \textbf{Región de rechazo} \\
    \midrule
    Bilateral (dos colas) & $\mu\neq\mu_0$ &
      Repartida en las dos colas, $\alpha/2$ en cada una \\
    Unilateral derecha & $\mu>\mu_0$ & Toda en la cola derecha, $\alpha$ \\
    Unilateral izquierda & $\mu<\mu_0$ & Toda en la cola izquierda, $\alpha$ \\
    \bottomrule
  \end{tabular}
  \end{center}

  \vspace{0.2cm}
  \begin{columns}[T,onlytextwidth]
    \begin{column}{0.48\textwidth}
      \begin{exampleblock}{Valores críticos con $\alpha=0{,}05$}
        \footnotesize
        \centering
        Bilateral: $z=\pm 1{,}96$ \\
        Unilateral: $z=1{,}645$ (o $-1{,}645$)
      \end{exampleblock}
    \end{column}
    \begin{column}{0.48\textwidth}
      \begin{alertblock}{Por qué importa}
        \footnotesize
        Con el mismo $\alpha$, la prueba unilateral rechaza más fácil: toda la
        región de rechazo está de un solo lado. Por eso el tipo debe decidirse
        \destacar{antes} de ver los datos.
      \end{alertblock}
    \end{column}
  \end{columns}
\end{frame}

\begin{frame}{Los dos errores posibles}
  \footnotesize
  \begin{center}
  \begin{tabular}{@{}lcc@{}}
    \toprule
    & \textbf{$H_0$ es verdadera} & \textbf{$H_0$ es falsa} \\
    \midrule
    \textbf{Se rechaza $H_0$} &
      \cellcolor{uniCrema}Error tipo I ($\alpha$) & Decisión correcta \\
    \textbf{No se rechaza $H_0$} &
      Decisión correcta & \cellcolor{uniCrema}Error tipo II ($\beta$) \\
    \bottomrule
  \end{tabular}
  \end{center}

  \vspace{0.25cm}
  \begin{columns}[T,onlytextwidth]
    \begin{column}{0.48\textwidth}
      \begin{alertblock}{Error tipo I}
        \footnotesize
        Condenar a un inocente: concluir que hay efecto cuando no lo hay. Su
        probabilidad es $\alpha$, y \azul{la fija el investigador}.
      \end{alertblock}
    \end{column}
    \begin{column}{0.48\textwidth}
      \begin{block}{Error tipo II}
        \footnotesize
        Absolver a un culpable: no detectar un efecto que sí existe. Su
        probabilidad es $\beta$. La \azul{potencia} de la prueba es
        $1-\beta$.
      \end{block}
    \end{column}
  \end{columns}

  \vspace{0.15cm}
  {\scriptsize\color{uniGris} Reducir $\alpha$ aumenta $\beta$ y viceversa. La
  única forma de reducir ambos a la vez es \destacar{aumentar $n$}.}
\end{frame}

\begin{frame}{Nivel de significancia y valor $p$}
  \begin{columns}[T,onlytextwidth]
    \begin{column}{0.48\textwidth}
      \begin{block}{Nivel de significancia $\alpha$}
        \footnotesize
        Probabilidad máxima de error tipo I que se está dispuesto a aceptar.
        Valores usuales: $0{,}10$, $0{,}05$ y $0{,}01$.

        \vspace{0.1cm}
        Se fija \azul{antes} de calcular nada.
      \end{block}
    \end{column}
    \begin{column}{0.48\textwidth}
      \begin{exampleblock}{Valor $p$}
        \footnotesize
        Probabilidad de obtener un resultado \azul{tan extremo o más} que el
        observado, \emph{suponiendo que $H_0$ es verdadera}.

        \vspace{0.1cm}
        Mide qué tan sorprendente es la muestra bajo $H_0$.
      \end{exampleblock}
    \end{column}
  \end{columns}

  \vspace{0.25cm}
  \begin{alertblock}{Regla de decisión}
    \footnotesize
    \centering
    \destacar{Si $p \le \alpha$, se rechaza $H_0$. Si $p > \alpha$, no se
    rechaza.}

    \vspace{0.1cm}
    {\scriptsize En palabras: si lo observado sería muy raro bajo $H_0$,
    dejamos de creer en $H_0$.}
  \end{alertblock}
\end{frame}

\begin{frame}{Lo que el valor $p$ NO es}
  \small
  \begin{alertblock}{Tres interpretaciones incorrectas}
    \footnotesize
    \begin{itemize}
      \item No es la probabilidad de que $H_0$ sea verdadera. $H_0$ no tiene
            probabilidad: es una afirmación fija.
      \item No es la probabilidad de haberse equivocado al rechazar.
      \item Un $p$ pequeño no significa que el efecto sea \azul{grande} ni
            \azul{importante}: solo que es difícil de explicar por azar. Con
            $n$ enorme, diferencias triviales dan $p$ minúsculos.
    \end{itemize}
  \end{alertblock}

  \vspace{0.2cm}
  \begin{block}{Significancia estadística $\neq$ relevancia práctica}
    \footnotesize
    Que un cambio de 0{,}3 minutos en el tiempo de atención sea
    ``estadísticamente significativo'' no implica que valga la pena invertir en
    él. Esa segunda decisión es del administrador, no del estadístico.
  \end{block}
\end{frame}

\begin{frame}{Ejercicio 1 — en clase}
  \begin{alertblock}{En parejas (9 minutos)}
    Plantee $H_0$ y $H_1$ e indique si la prueba es unilateral o bilateral.

    \vspace{0.25cm}
    \begin{enumerate}
      \item La norma dice que el contenido debe ser de 2{,}5 kg; se quiere
            verificar si la máquina está descalibrada.
      \item Se afirma que un nuevo proceso reduce el tiempo promedio, que era
            de 15 minutos.
      \item Se quiere demostrar que más del 60\% de los clientes está
            satisfecho.
      \item Se quiere verificar si la proporción de defectuosos cambió respecto
            al 5\% histórico.
    \end{enumerate}
  \end{alertblock}
\end{frame}

\begin{frame}{Ejercicio 1 — solución}
  \footnotesize
  \begin{center}
  \begin{tabular}{@{}clll@{}}
    \toprule
    \textbf{\#} & \textbf{$H_0$} & \textbf{$H_1$} & \textbf{Tipo} \\
    \midrule
    1 & $\mu=2{,}5$    & $\mu\neq 2{,}5$ & Bilateral \\
    2 & $\mu\ge 15$    & $\mu<15$        & Unilateral izquierda \\
    3 & $p\le 0{,}60$  & $p>0{,}60$      & Unilateral derecha \\
    4 & $p=0{,}05$     & $p\neq 0{,}05$  & Bilateral \\
    \bottomrule
  \end{tabular}
  \end{center}

  \vspace{0.25cm}
  \begin{exampleblock}{Cómo se reconoce el tipo}
    \footnotesize
    ``Descalibrada'' y ``cambió'' no dicen en qué dirección: \azul{bilateral}.
    ``Reduce'' y ``más del'' señalan una dirección: \azul{unilateral}, y el
    sentido de la desigualdad de $H_1$ indica cuál cola.
  \end{exampleblock}
\end{frame}

\begin{frame}{Ejercicio 2 — en clase}
  \begin{alertblock}{Individual (8 minutos)}
    Una prueba de hipótesis sobre el tiempo promedio de atención arrojó
    \destacar{$p=0{,}032$}.

    \vspace{0.25cm}
    \begin{enumerate}
      \item ¿Cuál es la decisión con $\alpha=0{,}05$?
      \item ¿Y con $\alpha=0{,}01$?
      \item Redacte la conclusión del punto 1 en lenguaje del negocio.
      \item ¿Es correcto decir que ``hay 3,2\% de probabilidad de que $H_0$
            sea verdadera''?
    \end{enumerate}
  \end{alertblock}
\end{frame}

\begin{frame}{Ejercicio 2 — solución}
  \small
  \begin{block}{Decisiones}
    \footnotesize
    \begin{itemize}
      \item Con $\alpha=0{,}05$: como $0{,}032\le 0{,}05$, \azul{se rechaza
            $H_0$}.
      \item Con $\alpha=0{,}01$: como $0{,}032>0{,}01$, \azul{no se rechaza
            $H_0$}.
    \end{itemize}
  \end{block}

  \vspace{0.15cm}
  \begin{exampleblock}{Punto 3}
    \footnotesize
    ``Con un nivel de significancia del 5\%, la evidencia muestral es
    suficiente para concluir que el tiempo promedio de atención difiere del
    valor de referencia''.
  \end{exampleblock}

  \vspace{0.1cm}
  \begin{alertblock}{Punto 4}
    \footnotesize
    \destacar{No.} El valor $p$ se calcula \emph{suponiendo} que $H_0$ es
    verdadera; no puede ser al mismo tiempo la probabilidad de que lo sea.
  \end{alertblock}
\end{frame}

\begin{frame}[fragile]{El error tipo I, visto en R}
  \scriptsize
  Si $H_0$ es \azul{verdadera} y aun así rechazamos el 5\% de las veces,
  ¿cuánto es ese 5\% en la práctica?
\begin{lstlisting}[style=Rmini]
set.seed(2026)
# H0 verdadera: la poblacion SI tiene media 100
rechazos <- replicate(2000, {
  x <- rnorm(25, mean = 100, sd = 15)
  t.test(x, mu = 100)$p.value < 0.05      # rechazo con alfa = 0.05
})
mean(rechazos)     # proporcion de veces que nos equivocamos
\end{lstlisting}
\begin{lstlisting}[style=Rout]
[1] 0.0505
\end{lstlisting}

  \vspace{-0.05cm}
  \begin{alertblock}{Lo que acaba de ocurrir}
    \scriptsize
    En 2\,000 estudios donde $H_0$ \destacar{era cierta}, la rechazamos 101
    veces: el 5\%, exactamente $\alpha$. \azul{El error tipo I no es un
    accidente: es el precio fijado de antemano.} Bajar $\alpha$ a 0,01 reduce
    esos falsos positivos, pero a cambio aumenta el error tipo II.
  \end{alertblock}

  \vspace{0.05cm}
  {\tiny\color{uniGris} Cambie \fx{mean = 100} por \fx{mean = 105} y la misma
  simulación le mostrará la \destacar{potencia} de la prueba.}
\end{frame}

\begin{frame}{Cierre de la sesión}
  \begin{block}{Lo que debe quedar claro hoy}
    \begin{itemize}
      \item $H_0$ lleva la igualdad y es el statu quo; $H_1$ es lo que se
            quiere probar.
      \item Bilateral si $H_1$ usa $\neq$; unilateral si usa $>$ o $<$.
      \item Error tipo I ($\alpha$): rechazar una $H_0$ verdadera. Tipo II
            ($\beta$): no rechazar una falsa.
      \item Regla: si $p\le\alpha$ se rechaza $H_0$.
      \item No rechazar no es demostrar que $H_0$ sea cierta.
    \end{itemize}
  \end{block}

  \vspace{0.2cm}
  \begin{alertblock}{Trabajo independiente para la próxima sesión}
    \footnotesize
    Resolver los \destacar{20 ejercicios de la lámina siguiente}. En la sesión
    14 aplicaremos los cinco pasos completos y usaremos software.
  \end{alertblock}
\end{frame}

\begin{frame}{Ejercicios propuestos}
  \vspace{-0.1cm}
  \begin{beamercolorbox}[wd=\textwidth,sep=0.10cm]{block body}
    \centering\scriptsize
    \textbf{1--8:} plantee $H_0$ y $H_1$ e indique el tipo. \quad
    \textbf{9--20:} conceptos y decisiones.
  \end{beamercolorbox}
  \listatiny
  \vspace{0.15cm}
  \begin{columns}[T,onlytextwidth]
    \begin{column}{0.47\textwidth}
      \begin{enumerate}\setlength{\itemsep}{0.16cm}
        \item ``El peso promedio es 500 g''
        \item ``El tiempo promedio bajó de 20 min''
        \item ``Más del 40\% usa la app''
        \item ``La proporción de defectuosos cambió''
        \item ``El gasto promedio supera 80 mil''
        \item ``La media es menor que 12''
        \item ``El proceso no alteró la media de 3{,}2''
        \item ``Menos del 10\% incumple''
        \item ¿Qué signo lleva siempre $H_0$?
        \item ¿Qué es el error tipo I?
      \end{enumerate}
    \end{column}
    \begin{column}{0.50\textwidth}
      \begin{enumerate}\setlength{\itemsep}{0.16cm}
        \setcounter{enumi}{10}
        \item ¿Qué es el error tipo II?
        \item ¿Quién fija $\alpha$?
        \item ¿Qué es la potencia de la prueba?
        \item $p=0{,}032$ y $\alpha=0{,}05$: decisión
        \item $p=0{,}032$ y $\alpha=0{,}01$: decisión
        \item $p=0{,}18$ y $\alpha=0{,}10$: decisión
        \item $z$ crítico bilateral con $\alpha=0{,}05$
        \item $z$ crítico unilateral con $\alpha=0{,}05$
        \item ¿$p$ es la probabilidad de que $H_0$ sea verdadera?
        \item ¿Cómo se reducen $\alpha$ y $\beta$ a la vez?
      \end{enumerate}
    \end{column}
  \end{columns}
\end{frame}

\begin{frame}{Respuestas}{Autoevaluación de los ejercicios propuestos}
  \listatiny
  \vspace{0.15cm}
  \begin{columns}[T,onlytextwidth]
    \begin{column}{0.47\textwidth}
      \begin{enumerate}\setlength{\itemsep}{0.16cm}
        \item $\mu=500$ vs. $\mu\neq 500$; bilateral
        \item $\mu\ge 20$ vs. $\mu<20$; unilateral izq.
        \item $p\le 0{,}40$ vs. $p>0{,}40$; unilateral der.
        \item $p=p_0$ vs. $p\neq p_0$; bilateral
        \item $\mu\le 80$ vs. $\mu>80$; unilateral der.
        \item $\mu\ge 12$ vs. $\mu<12$; unilateral izq.
        \item $\mu=3{,}2$ vs. $\mu\neq 3{,}2$; bilateral
        \item $p\ge 0{,}10$ vs. $p<0{,}10$; unilateral izq.
        \item La igualdad ($=$, $\le$ o $\ge$)
        \item Rechazar $H_0$ siendo verdadera
      \end{enumerate}
    \end{column}
    \begin{column}{0.50\textwidth}
      \begin{enumerate}\setlength{\itemsep}{0.16cm}
        \setcounter{enumi}{10}
        \item No rechazar $H_0$ siendo falsa
        \item El investigador, antes de calcular
        \item $1-\beta$
        \item Se rechaza $H_0$
        \item No se rechaza $H_0$
        \item No se rechaza $H_0$
        \item $\pm 1{,}96$
        \item $1{,}645$ (o $-1{,}645$)
        \item No
        \item Aumentando $n$
      \end{enumerate}
    \end{column}
  \end{columns}
\end{frame}

\begin{frame}{Referencias}
  \footnotesize
  \begin{itemize}
    \item[{[1]}] Díaz Rodríguez, M. (2022). \textit{Estadística inferencial
      aplicada}. Universidad del Norte.
      \texttt{https://elibro.net/es/lc/uniminuto/titulos/221614}
    \item[{[2]}] Llinás Solano, H. (2014). \textit{Introducción a la
      estadística con aplicaciones en ciencias sociales}. Universidad del
      Norte. \texttt{https://elibro.net/es/lc/uniminuto/titulos/70062}
    \item[{[3]}] Contento Rubio, M. R. (2019). \textit{Estadística con
      aplicaciones en R}. Editorial Utadeo.
      \texttt{https://elibro.net/es/lc/uniminuto/titulos/220926}
  \end{itemize}

  \vspace{0.35cm}
  \begin{center}
    {\Large\color{uniAzul}\textbf{¿Preguntas?}}
  \end{center}
\end{frame}

\end{document}
